\section{The composed relation maps}
\label{sec:composed-maps}

This appendix supplies the full-space formulas behind
Corollary~\ref{cor:oriented-relations} and
Proposition~\ref{prop:catalytic}; it is not a second proof of either.
Both proofs use Lemma~\ref{lem:composition}, which requires the maps
applied after the separator to be fixed linear maps on all target
variables, including those of terms of positive degree, whose first-leg
label \(g'\) differs from the sector label \(g\) of the other two legs.
For one relation we write out the separated polynomial with all of its
terms, \eqref{eq:composed-relation-polynomial}, and the reindexing on the
full target spaces, \eqref{eq:reindexing};
Section~\ref{subsec:composed-orientations} checks the later regroupings.

\subsection{The data of one relation}

We use the integer counts of Section~\ref{sec:rounding} at rounding scale
\(K\), the type scale \(N\), and the set \(\mathcal J=\mathcal
I_c\sqcup\mathcal I_t\) of active relations of
Section~\ref{sec:extraction}. The row of a relation \(i\in\mathcal J\) is
denoted \(r_i\).

\emph{Splits.} A split \(i=(r,x)\in\mathcal I_c\) has row \(r_i=r\), source
column \(b_i=r-1+x\), width \(h_i=h_{r,x}\) and auxiliary dimension
\(d_i=h_i+1\), and it is applied \(w_i=w_{r,x}\) times. A child \(y\) of
the split has column \(r-1+y=b_i+h_i-2j\) with \(j=(x+h_i-y)/2\), so it is
the diagonal block of Lemma~\ref{lem:multiway} with filtration index
\(j\). At type scale \(N\) the split uses the type
\((n^{(i)}_0,\ldots,n^{(i)}_{h_i})\) given by
\begin{equation}\label{eq:filtration-counts}
 n^{(i)}_j=N\,n_{r,x,y}\quad\text{with }y=x+h_i-2j
 \qquad(0\le j\le h_i),
\end{equation}
so that \(\sum_jn^{(i)}_j=Nw_i\) and the sector
count~\eqref{eq:sector-counts} is \(M_i=(Nw_i)!/\prod_jn^{(i)}_j!\). These
counts are indexed by filtration index, whereas the child counts
\(n_{i,y}=n_{r,x,y}\) of Section~\ref{sec:extraction} are indexed by child
state: for instance, \(n^{(i)}_0=Nn_{i,x+h_i}\) refers to the longest
child, of column \(b_i+h_i\), whereas \(n_{i,0}\) refers to the child at
state \(0\).

\emph{Triplings.} A tripling \(i\in\mathcal I_t\) is applied
\(z_{r_i}\) times. Its middle count is
\(m_{r_i}=\lfloor(Nz_{r_i}+1)/3\rfloor\), and its sector count \(M_i\) is
the count \(M_{r_i}\) of~\eqref{eq:sector-counts}, that is, the count
\(M_{\mathrm t}(Nz_{r_i},m_{r_i})\) of Lemma~\ref{lem:tripling}.

The \emph{source tensor} \(\mathcal U_i\) and the \emph{canonical sector
tensor} \(\mathcal V_i\) of relation \(i\) are
\[
 \mathcal U_i=
 \begin{cases}
 \bigl(C(r_i,b_i)B_X(d_i)\bigr)^{\otimes Nw_i},&i\in\mathcal I_c,\\[2pt]
 C(r_i,4r_i-1)^{\otimes Nz_{r_i}},&i\in\mathcal I_t,
 \end{cases}
 \qquad
 \mathcal V_i=
 \begin{cases}
 \prod_{j=0}^{h_i}
 C(r_i,b_i+h_i-2j)^{\otimes n^{(i)}_j},&i\in\mathcal I_c,\\[2pt]
 C(r_i,r_i)^{\otimes(Nz_{r_i}-m_{r_i})}
 \bigl(C(r_i,r_i)^{(23)}\bigr)^{\otimes m_{r_i}},&i\in\mathcal I_t.
 \end{cases}
\]
By Lemmas~\ref{lem:multiway} and~\ref{lem:tripling}, every sector of the
type restriction of \(\mathcal U_i\) is isomorphic to \(\mathcal V_i\).

\emph{Words and sectors.} Let \(\mathsf T_i\) be the type restriction of
Lemma~\ref{lem:multiway} or Lemma~\ref{lem:tripling}, which acts on the
second and third legs only, and index the tensor positions of
\(\mathcal U_i\) by \(\theta\). The \(M_i\) retained words, of type
\((n^{(i)}_j)_j\) for a split and with exactly \(m_{r_i}\) middle labels
for a tripling, are numbered \(g=0,\ldots,M_i-1\) in lexicographic order,
with letters ordered by filtration index for a split and as left, middle,
right for a tripling; the coordinates whose word has number \(g\) form
sector \(g\). By the proofs of the two lemmas, every term that survives the
type restriction has the same retained word on its second and third legs.
Hence
\begin{equation}\label{eq:type-restricted-relation}
 \mathsf T_i(\mathcal U_i)=\sum_{g=0}^{M_i-1}T_{i,g},\qquad
 T_{i,g}=\sum_{\alpha,\beta,\gamma}T_{i,g}[\alpha,\beta,\gamma]\,
 x_\alpha y_{g,\beta}z_{g,\gamma},
\end{equation}
where \(x_\alpha\) runs over the coordinates of the shared first leg
\(X_i\), and \(y_{g,\beta}\) and \(z_{g,\gamma}\) over those of sector
\(g\), which span \(Y_{i,g}\) and \(Z_{i,g}\); this is the form to which
Proposition~\ref{prop:separator} applies.

\emph{Coefficient arrays.} The indices are multi-indices
\(\alpha=(\alpha_\theta)_\theta\), \(\beta=(\beta_\theta)_\theta\) and
\(\gamma=(\gamma_\theta)_\theta\). For a split whose word \(g\) has letter
\(j_\theta\) at position \(\theta\), in the determinant bases and before
any positions are reordered, \(\alpha_\theta\) indexes a monomial of
\(V_{r_i}\), and \(\beta_\theta\) and \(\gamma_\theta\) index basis vectors
of filtration index \(j_\theta\) in the input and output spaces, with
ranges
\begin{equation}\label{eq:coefficient-ranges}
 0\le\alpha_\theta<r_i,\qquad
 0\le\beta_\theta<b_i+h_i-2j_\theta,\qquad
 0\le\gamma_\theta<r_i+b_i+h_i-2j_\theta-1,
\end{equation}
and the coefficients are
\begin{equation}\label{eq:coefficient-array}
 T_{i,g}[\alpha,\beta,\gamma]=\prod_\theta
 \mathbf 1_{\{\gamma_\theta=\alpha_\theta+\beta_\theta\}}.
\end{equation}
Indeed, in a position with letter \(j\) the type restriction keeps exactly
the diagonal block \(C(r_i,b_i+h_i-2j)\), and on this block multiplication
by the monomial with index \(\alpha\) sends monomial index \(\beta\) to
\(\beta+\alpha\) with coefficient one
(Lemma~\ref{lem:determinant-basis}(b),(c) and the computation that
follows it in Section~\ref{subsec:splits}). For a tripling,
\eqref{eq:coefficient-array} holds with \(0\le\alpha_\theta<r_i\) when
\(\beta_\theta\) and \(\gamma_\theta\) are the original coordinates of
\(C(r_i,4r_i-1)\), restricted to the two intervals of the letter at
position \(\theta\). The shifts of these interval coordinates and the
reversal of the first index at middle positions are part of the
reindexing \(\mathsf R_i\) below.

\subsection{The separated polynomial of one relation}

Put \(L_i=2M_i-1\) and \(\zeta=\zeta_{L_i}\), and let \(\mathsf S_i(t)\)
be the substitutions of Proposition~\ref{prop:separator} for the \(M_i\)
sectors in~\eqref{eq:type-restricted-relation}, with the phases of Step~1
and the weights of Step~2 of its proof combined. We suppress the index
\(i\) on the new variables \(X_{\alpha,g'}\), \(Y_{g,\beta,e}\) and
\(Z_{g,\gamma,e'}\), where \(g',e,e'\in\{0,\ldots,M_i-1\}\). In source copy
\(\ell\), \(0\le\ell\le L_i-1\), the substitutions are
\[
 x_\alpha\mapsto\sum_{g'}\zeta^{\ell g'}t^{g'^2}X_{\alpha,g'},\quad
 y_{g,\beta}\mapsto\sum_e\zeta^{\ell(e-g)}t^{2ge-g^2}Y_{g,\beta,e},\quad
 z_{g,\gamma}\mapsto\frac1{L_i}\sum_{e'}\zeta^{-\ell e'}t^{-2ge'}
 Z_{g,\gamma,e'}.
\]
By Steps~1 and~2 of the proof of Proposition~\ref{prop:separator},
\(\mathsf S_i(t)\), applied after the type restrictions in the \(L_i\)
source copies, transforms \(L_i\mathcal U_i\) into exactly
\begin{equation}\label{eq:composed-relation-polynomial}
 \widetilde T_i(t)=
 \sum_{\substack{0\le g,g',e,e'<M_i\\e-e'+g'-g=0}}
 t^{(g'-g)^2}
 \sum_{\alpha,\beta,\gamma}
 T_{i,g}[\alpha,\beta,\gamma]\,
 X_{\alpha,g'}Y_{g,\beta,e}Z_{g,\gamma,e'},
\end{equation}
where \(\alpha,\beta,\gamma\) run over the ranges of word \(g\); with
\eqref{eq:coefficient-array}, every coefficient is explicit, including
those of the terms of positive degree. The exponents are
\(\delta^2\) with \(\delta=|g'-g|\le M_i-1\), and we write
\(\widetilde T_i(t)=\sum_{\delta=0}^{M_i-1}t^{\delta^2}
\widetilde T_{i,\delta^2}\), where the coefficient
\(\widetilde T_{i,\delta^2}\) of \(t^{\delta^2}\) collects the terms with
\(|g'-g|=\delta\). The constant term consists of
the terms with \(g'=g\), for which \(e=e'\), so
\(\widetilde T_{i,0}=\bigoplus_{g=0}^{M_i-1}T_{i,g}B_X(M_i)\), with
variables \(X_{\alpha,g}\), \(Y_{g,\beta,e}\) and \(Z_{g,\gamma,e}\).

\subsection{Reindexing on the full target spaces}

The target spaces of the separator are
\(\widetilde X_i=\bigoplus_{g'}X_i^{(g')}\), where \(X_i^{(g')}\) is the
copy of \(X_i\) with coordinates \(X_{\alpha,g'}\), and
\(\widetilde Y_i=\bigoplus_g(Y_{i,g}\otimes\C^{M_i})\) and
\(\widetilde Z_i=\bigoplus_g(Z_{i,g}\otimes\C^{M_i})\), where the factor
\(\C^{M_i}\) carries the index \(e\), respectively \(e'\).

For each sector \(g\), let \(\mathsf R^X_{i,g}\), \(\mathsf R^Y_{i,g}\)
and \(\mathsf R^Z_{i,g}\) be the bijections of coordinates from \(X_i\),
\(Y_{i,g}\) and \(Z_{i,g}\) to the three legs of \(\mathcal V_i\) that
identify \(T_{i,g}\) with \(\mathcal V_i\). For a split, they reorder the
positions by increasing filtration index in word \(g\), keeping the
original order of positions with equal filtration index; after this
reordering, \eqref{eq:coefficient-array} with the
ranges~\eqref{eq:coefficient-ranges} is the coefficient array of
\(\mathcal V_i\). For a tripling, they put all outer positions before all
middle positions, again keeping the original order within each group;
they subtract from \(\beta_\theta\) and \(\gamma_\theta\) the left
endpoints of their intervals, and at middle positions they also reverse
the first index, \(\alpha_\theta\mapsto r_i-1-\alpha_\theta\). By the
proof of Lemma~\ref{lem:tripling}, this identifies every outer block with
\(C(r_i,r_i)\) and every middle block with \(C(r_i,r_i)^{(23)}\); left and
right blocks give the same canonical block. Define the reindexing on the
full target spaces by
\begin{equation}\label{eq:reindexing}
 \mathsf R_i=
 \Bigl(
  \bigoplus_{g'}\mathsf R^X_{i,g'},\
  \bigoplus_g\bigl(\mathsf R^Y_{i,g}\otimes\mathrm{id}_{\C^{M_i}}\bigr),\
  \bigoplus_g\bigl(\mathsf R^Z_{i,g}\otimes\mathrm{id}_{\C^{M_i}}\bigr)
 \Bigr).
\end{equation}
The first component reorders the copy \(X_i^{(g')}\) according to word
\(g'\), and the other two reorder sector \(g\) according to word \(g\). A
monomial of positive degree has \(g'\ne g\), so its legs are reordered
according to different words and its image is not in a copy of
\(\mathcal V_i\); exactification removes such terms later. Nevertheless
\(\mathsf R_i\) is a bijection between the coordinates of the full spaces
and does not depend on \(t\), which is what Lemma~\ref{lem:composition}
requires.

For an orientation \(\pi\in\mathcal A_{\sigma_{r_i}}\), where
\(\sigma_{r_i}=(-1)^{a-r_i}\) is the sign of row \(r_i\), the complete map
of relation \(i\) is
\begin{equation}\label{eq:complete-relation-map}
 \Psi_{i,\pi}(t)=\mathsf R_i^\pi\circ\mathsf S_i(t)^\pi\circ
 \Bigl(\bigoplus_{\ell=0}^{L_i-1}\mathsf T_i^\pi\Bigr).
\end{equation}
It acts on \(L_i\mathcal U_i^\pi\), and its separator uses as shared leg
the global leg \(\pi(1)\), which contains leg~1 of \(\mathcal U_i\). By
Lemma~\ref{lem:composition}, \(\mathsf R_i^\pi\) acts on
\(\widetilde T_i(t)^\pi\) coefficient by coefficient. In the constant term
every monomial has \(g'=g\), so \(\mathsf R_i\) reindexes block \(g\)
according to the single word \(g\) and maps \(T_{i,g}B_X(M_i)\) onto
\(\mathcal V_iB_X(M_i)\). Hence
\begin{equation}\label{eq:relation-constant-term}
 \Psi_{i,\pi}(t)(L_i\mathcal U_i^\pi)
 =\bigoplus_{g=0}^{M_i-1}\bigl(\mathcal V_iB_X(M_i)\bigr)^\pi
 +\sum_{\delta=1}^{M_i-1}t^{\delta^2}
   \bigl(\mathsf R_i(\widetilde T_{i,\delta^2})\bigr)^\pi,
\end{equation}
and the first term is a full direct sum on all three legs.

\begin{example}[Two tensor positions]\label{ex:two-positions}
Let \(i\) be a split of width \(h_i=1\) with type
\(n^{(i)}_0=n^{(i)}_1=1\) (for instance \(N=1\) and \(w_i=2\), each child
used once), and write \(r=r_i\) and \(b=b_i\ge2\). Then
\(\mathcal U_i=(C(r,b)B_X(2))^{\otimes2}\),
\(\mathcal V_i=C(r,b+1)C(r,b-1)\), and \(M_i=2\): sector \(0\) has the
word \((0,1)\) and sector \(1\) the word \((1,0)\). By
\eqref{eq:coefficient-ranges}, the array of \(T_{i,0}\) has ranges
\(\beta_1<b+1\), \(\gamma_1<r+b\), \(\beta_2<b-1\), \(\gamma_2<r+b-2\), so
it is the array of \(\mathcal V_i\), and the array of \(T_{i,1}\) is the
same with the two positions exchanged. Hence
\(\mathsf R^X_{i,0},\mathsf R^Y_{i,0},\mathsf R^Z_{i,0}\) are identities,
and \(\mathsf R^X_{i,1},\mathsf R^Y_{i,1},\mathsf R^Z_{i,1}\) exchange the
two positions.

The separator uses \(L_i=3\) copies, and \(g',e,e'\in\{0,1\}\). The
condition \(e-e'+g'-g=0\) leaves six index tuples \((g,g',e,e')\): the
four with \(g'=g\) and \(e=e'\), which give
\(T_{i,0}B_X(2)\oplus T_{i,1}B_X(2)\), mapped by \(\mathsf R_i\) onto
\(2\,\mathcal V_iB_X(2)\); and \((0,1,0,1)\) and \((1,0,1,0)\), which give
the terms of degree one
\[
 t\sum_{\alpha,\beta,\gamma}T_{i,0}[\alpha,\beta,\gamma]\,
   X_{\alpha,1}Y_{0,\beta,0}Z_{0,\gamma,1}
 +t\sum_{\alpha,\beta,\gamma}T_{i,1}[\alpha,\beta,\gamma]\,
   X_{\alpha,0}Y_{1,\beta,1}Z_{1,\gamma,0}
\]
of~\eqref{eq:composed-relation-polynomial}. In the first of them the first
leg has label \(g'=1\) and the other legs have label \(g=0\), so
\(\mathsf R_i\) exchanges the two positions of \(\alpha\) but not those of
\(\beta\) and \(\gamma\). The image pairs the first-leg factor of each
position with the second- and third-leg factors of the other position; it
is a well-defined term of degree one, but not a copy of \(\mathcal V_i\),
and exactification removes it.
\end{example}

\subsection{Three orientations and all relations}
\label{subsec:composed-orientations}

Tensoring the maps \(\Psi_{i,\pi}(t)\) over \(\pi\in\mathcal A_\sigma\),
where \(\sigma=\sigma_{r_i}\), makes every global leg of the target a
direct sum of blocks indexed by triples \((g_\pi)_{\pi\in\mathcal
A_\sigma}\) of independently chosen labels.
By~\eqref{eq:relation-constant-term}, the constant term has \(M_i^3\) such
blocks, each carrying \(\bigotimes_\pi(\mathcal V_iB_X(M_i))^\pi\). The
proof of Corollary~\ref{cor:oriented-relations} identifies the source
\(\bigotimes_\pi\mathcal U_i^\pi\) and each of these blocks with the two
sides of~\eqref{eq:split-relation} or~\eqref{eq:tripling-relation}: it
reorders tensor factors within each global leg, combines the three
oriented dot products by~\eqref{eq:coset-dot-products}, and, for a
tripling, uses~\eqref{eq:cyclic-tripling-balance}, including the diagonal
isomorphism \(S^+_{r,r}\cong S^-_{r,r}\) of~\eqref{eq:cyclic-identities}.
We apply each of these identifications as the same fixed coordinate map
on every block of a global leg. It is therefore defined also on monomials
of positive degree whose labels \(g_\pi\) differ, and it does not depend
on \(t\). With Lemma~\ref{lem:composition}, each relation thus gives the
degeneration of Corollary~\ref{cor:oriented-relations}, of degree at most
\(3(M_i-1)^2\), with every map defined on the full variable spaces.

The proof of Proposition~\ref{prop:catalytic} then tensors these maps over
all \(i\in\mathcal J\) with one parameter \(t\), and identifies factors
by the coordinate isomorphisms of~\eqref{eq:cyclic-identities} (at a row
boundary, \(S^\sigma_{r,b}\cong S^{-\sigma}_{b,r}\)) and
of~\eqref{eq:matrix-products}, again applied identically on every block
after all separators. Hence Lemma~\ref{lem:composition} applies to the
complete composed map, which is the
degeneration~\eqref{eq:catalytic-degeneration} of degree at most \(D\).
